\documentclass[11pt,a4paper]{article}
\usepackage[utf8]{inputenc}
\usepackage[T1]{fontenc}
\usepackage[english]{babel}
\usepackage{amsmath,amssymb,amsthm}
\\usepackage[a4paper,margin=0.65in]{geometry}
\emergencystretch=3em
\tolerance=600
\usepackage{graphicx}
\usepackage{array}
\usepackage{fancyhdr}
\usepackage[protrusion=true,expansion=false]{microtype}
\usepackage{hyperref}

% Fallbacks for non-standard commands
\providecommand{\half}{\tfrac{1}{2}}
\providecommand{\third}{\tfrac{1}{3}}

\pagestyle{fancy}
\fancyhf{}
\rhead{Cayley --- Collected Papers, Vol.~IX}
\lhead{Pages 504--528}
\cfoot{\thepage}

\title{\textbf{Arthur Cayley}\\\vspace{0.5em}
\Large Collected Mathematical Papers\\\vspace{0.3em}
\large Volume IX, Pages 504--528 (Papers 613--617)}
\author{Transcribed from the Original Publications}
\date{}

\begin{document}
\maketitle
\thispagestyle{empty}

\tableofcontents
\newpage

%% ========================================================================
%% PAPER [613] -- On the Group of Points G_4^1 on a Sextic Curve with Five Double Points
%%                Printed pages 504--507 (PDF pp. 526--529)
%% ========================================================================
\section*{\textbf{613.} On the Group of Points $G_{4}^{1}$ on a Sextic Curve with Five Double Points}
\addcontentsline{toc}{section}{\textbf{613.} On the Group of Points $G_{4}^{1}$ on a Sextic Curve with Five Double Points (pp.~504--507)}

\textit{[From the Mathematische Annalen, vol.~viii. (1875), pp.~359--362.]}

\bigskip

\noindent\textbf{[p.~504]}\quad The present note relates to a special group of points considered incidentally by MM.~Brill and N\"other in their paper ``Ueber die algebraischen Functionen und ihre Anwendung in der Geometrie,'' \textit{Math. Annalen}, t.~VII. pp.~268--310 (1874).

I recall some of the fundamental notions. We have a basis-curve which to fix the ideas may be taken to be of the order $n$, $= p + 1$, $\tfrac{1}{2}p(p - 3)$~dps, and therefore of the ``Geschlecht'' or deficiency~$p$; any curve of the order $n - 3$, $= p - 2$ passing through the $\tfrac{1}{2}p(p - 3)$~dps is said to be an \textit{adjoint curve}. We may have, on the basis-curve, a special group $G_{Q}^{q}$ of $Q$ points $(Q \geq 2p - 2)$; viz.\ this is the case when the $Q$ points are such that every adjoint curve through $Q - q$ of them---that is, every curve of the order $p - 2$ through $\tfrac{1}{2}p(p - 3)$~dps and the $Q - q$ points---passes through the remaining $q$ points of the group: the number $q$ may be termed the ``speciality'' of the group: if $q = 0$, the group is an ordinary one.

It may be observed that a special group $G_{Q}^{q}$ is chiefly noteworthy in the case where $Q - q$ is so small that the adjoint curve is not completely determined: thus if $p = 5$, viz.\ if the basis-curve be a sextic with 5~dps, then we may have a special group $G_{6}^{1}$, but there is nothing remarkable in this; the 6~points are intersections with the sextic of an arbitrary cubic through the 5~dps---the cubic of course intersects the sextic in the 5~dps counting as 10~points, and in 8 other points---and such cubic is completely determined by means of the 5~dps and any 4 of the 6~points. But contrariwise, there is something remarkable in the group $G_{4}^{1}$ about to be considered: viz.\ we have here on the sextic 4 points, such that every cubic through the 5~dps and through 3 of the 4 points (through 8 points in all) passes through the remaining one of the 4 points.

The whole number of intersections of the basis-curve with an adjoint, exclusive of the dps counting as $p(p - 3)$ points, is of course $= 2p - 2$: hence an adjoint through the $Q$ points of a group $G_{Q}^{q}$ meets the basis-curve besides in $R$, $= 2p - 2 - Q$, \noindent\textbf{[p.~505]}\quad points; we have then the ``Riemann-Roch'' theorem that these $R$ points form a special group $G_{R}^{r}$, where
\[
Q + R = 2p - 2,
\]
as just mentioned, and
\[
Q - R = 2q - 2r;
\]
viz.\ dividing in any manner the $2p - 2$ intersections of the basis-curve by an adjoint into groups of $Q$ and $R$ points respectively, these will be special groups, or at least one of them will be a special group, $G_{Q}^{q}$, $G_{R}^{r}$, such that their specialities $q$, $r$ are connected by the foregoing relation $Q - R = 2q - 2r$.

The Authors give (l.c., p.~293) a Table showing for a given basis-curve, or given value of $p$, and for a given value of $r$, the least value of $R$ and the corresponding values of $q$, $Q$: this table is conveniently expressed in the following form.

The least value of
\[
R = p - \frac{p}{r + 1} + r,
\]
and then
\[
q = \frac{p}{r + 1} - 1,
\]
\[
Q = p + \frac{p}{r + 1} - r - 2,
\]
where $\dfrac{p}{r + 1}$ denotes the integer equal to or next less than the fraction.

It is, I think, worth while to present the table in the more developed form:
\[
\begin{array}{|c|c|c|c|c|c|c|c|c|}
\hline
n & p & \text{Dps} & \multicolumn{6}{c|}{r =} \\
\cline{4-9}
 & & & 1 & 2 & 3 & 4 & 5 & 6 \\
\hline
4 & 3 & 0 & G_{2}^{1} & G_{3}^{2} & \cdot & \cdot & \cdot & \cdot \\
 & & & G_{4}^{2} & G_{6}^{3} & \cdot & \cdot & \cdot & \cdot \\
\hline
5 & 4 & 2 & G_{3}^{1} & G_{4}^{2} & G_{6}^{3} & \cdot & \cdot & \cdot \\
 & & & G_{4}^{2} & G_{6}^{3} & G_{8}^{4} & \cdot & \cdot & \cdot \\
\hline
6 & 5 & 5 & G_{4}^{1} & G_{4}^{2} & G_{6}^{3} & G_{8}^{4} & \cdot & \cdot \\
 & & & G_{4}^{2} & G_{6}^{3} & G_{8}^{4} & G_{10}^{5} & \cdot & \cdot \\
\hline
7 & 6 & 9 & G_{4}^{1} & G_{6}^{2} & G_{6}^{3} & G_{8}^{4} & G_{10}^{5} & \cdot \\
 & & & G_{6}^{2} & G_{6}^{3} & G_{8}^{4} & G_{10}^{5} & G_{12}^{6} & \cdot \\
\hline
8 & 7 & 14 & G_{4}^{1} & G_{6}^{2} & G_{8}^{3} & G_{10}^{4} & G_{10}^{5} & G_{12}^{6} \\
 & & & G_{6}^{2} & G_{8}^{3} & G_{8}^{4} & G_{10}^{5} & G_{12}^{6} & G_{14}^{7} \\
\hline
\end{array}
\]
where the table shows the values of $\dfrac{G_{R}^{r}}{G_{Q}^{q}}$ for any given values of $p$, $r$.

\noindent\textbf{[p.~506]}\quad I recur to the case $p = 5$ and the group $G_{4}^{1}$, which is the subject of the present note: viz.\ we have here a sextic curve with 5~dps, and on it a group of 4 points $G_{4}^{1}$, such that every cubic through the 5~dps and through 3 points of the group, 8 points in all, passes through the remaining 1 point.

MM.~Brill and N\"other show (by consideration of a rational transformation of the whole figure) that, given 2 points of the group, it is possible, and possible in 5 different ways, to determine the remaining 2 points of the group.

I remark that the 5~dps and the 4 points of the group form ``an ennead'' or system of the nine intersections of two cubic curves: and that the question is, given the 5~dps and 2 points on the sextic, to show how to determine on the sextic a pair of points forming with the 7 points an ennead: and to show that the number of solutions is $= 5$.

We have the following ``Geiser-Cotterill'' theorem:

If seven of the points of an ennead are fixed, and the eighth point describes a curve of the order $n$ passing $\alpha_{1}, \alpha_{2}, \ldots, \alpha_{7}$ times through the seven points respectively, then will the ninth point describe a curve of the order $\nu$ passing $a_{1}, a_{2}, \ldots, a_{7}$ times through the seven points respectively: where
\[
\nu = 8n - 3\Sigma\alpha,
\]
\[
a_{1} = 3n - \alpha_{1} - \Sigma\alpha,
\]
\[
\vdots
\]
\[
a_{7} = 3n - \alpha_{7} - \Sigma\alpha,
\]
and conversely
\[
n = 8\nu - 3\Sigma a,
\]
\[
\alpha_{1} = 3\nu - a_{1} - \Sigma a,
\]
\[
\vdots
\]
\[
\alpha_{7} = 3\nu - a_{7} - \Sigma a.
\]
(Geiser, \textit{Crelle-Borchardt}, t.~LXVII. (1867), pp.~78--90; the complete form, as just stated, and which was obtained by Mr~Cotterill, has not I believe been published); and also Geiser's theorem ``the locus of the coincident eighth and ninth points is a sextic passing twice through each of the seven points.''

The sextic and the curve $n$ intersect in $6n$ points, among which are included the seven points counting as $2\Sigma\alpha$ points: the number of the remaining points is $= 6n - 2\Sigma\alpha$. Similarly, the sextic and the curve $\nu$ intersect in $6\nu$ points, among which are included the seven points counting as $2\Sigma a$ points: the number of the remaining points is $6\nu - 2\Sigma a$ $(= 6n - 2\Sigma\alpha)$. The points in question are, it is clear, common intersections of the sextic, and the curves $n$, $\nu$: viz.\ of the intersections of the curves $n$, $\nu$, a number $6n - 2\Sigma\alpha$, $= 6\nu - 2\Sigma a$, $= 3n + 3\nu - 2\Sigma a - 2\Sigma\alpha$ lie on the sextic.

The curves $n$, $\nu$ intersect in $n\nu$ points, among which are included the seven points counting $\Sigma a\alpha$ times: the number of the remaining intersections is therefore \noindent\textbf{[p.~507]}\quad $n\nu - \Sigma a\alpha$, but among these are included the $3n + 3\nu - 2\Sigma a - 2\Sigma\alpha$ points on the sextic; omitting these, there remain $n\nu - 3(n + \nu) - \Sigma a\alpha + 2\Sigma a + 2\Sigma\alpha$ points, or, what is the same thing, $(n - 3)(\nu - 3) - \Sigma(a - 1)(\alpha - 1) - 2$ points: it is clear that these must form pairs such that, the eighth point being either point of a pair, the ninth point will be the remaining point of the pair: the number of pairs is of course
\[
\tfrac{1}{2}[(n - 3)(\nu - 3) - \Sigma(a - 1)(\alpha - 1) - 2],
\]
and we have thus the solution of the question, given the seven points to determine the number of pairs of points on the curve $n$ (or on the curve $\nu$) such that each pair may form with the seven points an ennead.

In particular, if $n = 6$; $\alpha_{1}, \alpha_{2}, \alpha_{3}, \alpha_{4}, \alpha_{5}, \alpha_{6}, \alpha_{7} = 2, 2, 2, 2, 2, 1, 1$ respectively, viz.\ if the curve be a sextic having 5 of the points for dps, and the remaining two for simple points, then we find $\nu = 12$; $a_{1}, a_{2}, a_{3}, a_{4}, a_{5}, a_{6}, a_{7} = 4, 4, 4, 4, 4, 5, 5$ respectively, and the number of pairs is
\[
= \tfrac{1}{2}[3 \cdot 9 - 5(2 - 1)(4 - 1) - 2], = \tfrac{1}{2}(27 - 15 - 2), = 5,
\]
viz.\ starting with the 5~dps and any 2 points of the group $G_{4}^{1}$ we can, in 5 different ways, determine the remaining 2 points of the group.

In reference to the number $3p - 3$ of parameters in the curves belonging to a given value of $p$, it may be remarked as follows. Such a curve is rationally transformable into a curve of the order $p + 1$ with $\tfrac{1}{2}p(p - 3)$ dps, and therefore containing $\tfrac{1}{2}(p + 1)(p + 4)$, $- \tfrac{1}{2}p(p - 3)$, $= 4p + 2$ parameters. Employing an arbitrary homographic transformation to establish any assumed relations between the parameters, the number is diminished to $4p + 2 - 8$, $= 4p - 6$; and again employing a rational transformation by means of adjoint curves of the order $p - 2$ drawn through the dps and $p - 3$ points of the curve---thereby transforming the curve into one of the same order $p + 1$ and deficiency $p$---then, assuming that the $p - 3$ parameters (or constants on which depend the positions of the $p - 3$ points) can be disposed of so as to establish $p - 3$ relations between the parameters and so further diminish the number by $p - 3$, the required number of parameters will finally be $4p - 6 - (p - 3) = 3p - 3$.

\medskip
\noindent\textit{Cambridge, 26th October, 1874.}

\newpage

%% ========================================================================
%% PAPER [614] -- On a Problem of Projection
%%                Printed pages 508--518 (PDF pp. 530--540)
%% ========================================================================
\section*{\textbf{614.} On a Problem of Projection}
\addcontentsline{toc}{section}{\textbf{614.} On a Problem of Projection (pp.~508--518)}

\textit{[From the Quarterly Journal of Pure and Applied Mathematics, vol.~xiii. (1875), pp.~19--29.]}

\bigskip

\noindent\textbf{[p.~508]}\quad I measure off on three rectangular axes the distances $\Omega X = \Omega Y = \Omega Z$, $= \theta$; and then, in a plane through $\Omega$ drawing in arbitrary directions the three lines $\Omega A$, $\Omega B$, $\Omega C$, $= a$, $b$, $c$ respectively, I assume that $A$, $B$, $C$ (fig.~1) are the parallel projections of $X$, $Y$, $Z$ respectively; viz.\ taking $\Omega O$ as the direction of the projecting lines, then $\Omega A$, $\Omega B$, $\Omega C$ being given in position and magnitude, we have to find $\theta$, and the position of the line $\Omega O$.

\begin{center}
\textit{[Fig.~1: Sphere with centre $\Omega$, showing axes $X$, $Y$, $Z$ and projections $A$, $B$, $C$ on a great circle, with projecting point $O$.]}
\end{center}

This is in fact a case of a more general problem solved by Prof.~Pohlke in 1853, (see the paper by Schwarz, ``Elementarer Beweis des Pohlke'schen Fundamentalsatzes der Axonometrie,'' \textit{Crelle}, t.~LXIII. (1864), pp.~309--314), viz.\ the three lines $\Omega X$, $\Omega Y$, $\Omega Z$ may be any three axes given in magnitude and direction, and their parallel projection \noindent\textbf{[p.~509]}\quad is to be similar to the three lines $\Omega A$, $\Omega B$, $\Omega C$. Schwarz obtains a very elegant construction, which I will first reproduce. We may imagine through $\Omega$ a plane cutting at right angles the projecting lines, say in the points $X'$, $Y'$, $Z'$: we have then in plano a triad of lines $\Omega X'$, $\Omega Y'$, $\Omega Z'$ which are an orthogonal projection of $\Omega X$, $\Omega Y$, $\Omega Z$; and are also an orthogonal projection of a plane triad similar to $\Omega A$, $\Omega B$, $\Omega C$; \textit{qu\^a} such last-mentioned projection, the triangles $\Omega Y'Z'$, $\Omega Z'X'$, $\Omega X'Y'$, must be proportional to the triangles $\Omega BC$, $\Omega CA$, $\Omega AB$; that is, we have to find an orthogonal projection of $\Omega X$, $\Omega Y$, $\Omega Z$, such that the triangles $\Omega Y'Z'$, $\Omega Z'X'$, $\Omega X'Y'$, which are the projections of $\Omega YZ$, $\Omega ZX$, $\Omega XY$ respectively, shall be in given ratios. There is no difficulty in the solution of this problem; referring everything to a sphere centre $\Omega$, let the normals to the planes $\Omega YZ$, $\Omega ZX$, $\Omega XY$, meet the sphere in the points $X''$, $Y''$, $Z''$ respectively, and the projecting line through $\Omega$ meet the sphere in the point $O$, then the projection of $\Omega YZ$ is to $\Omega YZ$ as $\cos OX'' : 1$; and the like as to the projections of $\Omega ZX$ and $\Omega XY$: that is, in the given spherical triangle $X''Y''Z''$, we have to find a point $O$, such that the cosines of the distances $OX''$, $OY''$, $OZ''$ are in given ratios: we have at once, through $X''$, $Y''$, $Z''$ respectively, three arcs meeting in the required point $O$.

The projecting lines being thus obtained, say these are the three parallel lines $X'$, $Y'$, $Z'$, we have next to draw through $\Omega$ a plane meeting these in the points $A'$, $B'$, $C'$ such that the triangle $A'B'C'$ is similar to the given triangle $ABC$; for this being so, the triangles $\Omega B'C'$, $\Omega C'A'$, $\Omega A'B'$ being the projections of, and therefore proportional to $\Omega Y'Z'$, $\Omega Z'X'$, $\Omega X'Y'$, that is, proportional to $\Omega BC$, $\Omega CA$, $\Omega AB$, will, it is clear, be similar to these triangles respectively; that is, we have the triad $\Omega A'$, $\Omega B'$, $\Omega C'$, a projection of $\Omega X$, $\Omega Y$, $\Omega Z$, and similar to the triad $\Omega A$, $\Omega B$, $\Omega C$, which is what was required.

It remains only to show how the given three parallel lines $X'$, $Y'$, $Z'$, not in the same plane, can be cut by a plane in a triangle similar to a given triangle $ABC$.

\begin{center}
\textit{[Fig.~2: Three parallel lines perpendicular to the plane of the paper through points $X$, $Y$, $Z$; auxiliary points $A''$, $D$, $E$, $K$, $B'$, $C'$ used in the construction.]}
\end{center}

Imagine the three lines at right angles to the plane of the paper, meeting the plane of the paper in the given points $X$, $Y$, $Z$ (fig.~2) respectively. On the base \noindent\textbf{[p.~510]}\quad $YZ$ describe a triangle $A''YZ$ similar to the given triangle $ABC$; and through $A''$, $X$ with centre on the line $YZ$, describe a circle meeting this line in the points $D$ and $E$. Then in the plane, through $YZ$ at right angles to the plane of the paper, we may draw a line meeting the lines $Y$, $Z$ in the points $B''$, $C''$ respectively, such that joining $XB''$, $XC''$ we obtain a triangle $XB''C''$ similar to $A''YZ$, that is, to the given triangle $ABC$.

Taking $K$ the centre of the circle, suppose that its radius is $= l$, and that we have $KY = \beta$, $KZ = \gamma$; also $YZ = \sigma$, $ZX = \tau$; $YA'' = \sigma''$, $ZA'' = \tau''$. If for a moment $x$, $y$ denote the coordinates of $A''$, then
\[
\tau^{2} = (x - \gamma)^{2} + y^{2}, = l^{2} + \gamma^{2} - 2\gamma x,
\]
and thence
\[
\gamma\sigma^{2} - \beta\tau^{2} = \gamma(l^{2} + \beta^{2}) - \beta(l^{2} + \gamma^{2}),
\]
that is,
\[
\gamma\sigma^{2} - \beta\tau^{2} = (\gamma - \beta)(l^{2} - \beta\gamma);
\]
viz.\ this is the equation of the circle in terms of the vectors $\sigma$, $\tau$; we have therefore in like manner
\[
\gamma\sigma''^{2} - \beta\tau''^{2} = (\gamma - \beta)(l^{2} - \beta\gamma).
\]

We may determine $\theta$ so as to satisfy the two equations
\[
\sigma''^{2} = \sigma^{2}\cos^{2}\theta + (l + \beta)^{2}\sin^{2}\theta,
\]
\[
\tau''^{2} = \tau^{2}\cos^{2}\theta + (l + \gamma)^{2}\sin^{2}\theta;
\]
in fact, these equations give
\[
\gamma\sigma''^{2} - \beta\tau''^{2} = (\gamma\sigma^{2} - \beta\tau^{2})\cos^{2}\theta + \{\gamma(l + \beta)^{2} - \beta(l + \gamma)^{2}\}\sin^{2}\theta,
\]
which, the left-hand side and the coefficients of $\cos^{2}\theta$ and $\sin^{2}\theta$ on the right-hand side being each $= (\gamma - \beta)(l^{2} - \beta\gamma)$, is, in fact, an identity.

But in the figure, if $\theta$, determined as above, denote the angle at $B$, then
\[
(XB'')^{2} = XY^{2} + YB''^{2} = \sigma^{2} + (l + \beta)^{2}\tan^{2}\theta,
\]
\[
(XC'')^{2} = XZ^{2} + ZC''^{2} = \tau^{2} + (l + \gamma)^{2}\tan^{2}\theta,
\]
that is,
\[
XB'' = \sigma''\sec\theta,\quad XC'' = \tau''\sec\theta,
\]
or, since $B''C'' = YZ\sec\theta\ [= (\gamma - \beta)\sec\theta]$, the triangle $XB''C''$ is, as mentioned, similar to the triangle $A''YZ$.

I was not acquainted with the foregoing construction when my paper was written; but the analytical investigation of the particular case is nevertheless interesting, and I proceed to consider it.

\noindent\textbf{[p.~511]}\quad Taking (fig.~1) $\Omega$ as the centre of a sphere and projecting on this sphere, we have $A$, $B$, $C$ given points on a great circle; and we have to find the point $O$, such that there may be a trirectangular triangle $XYZ$, the vertices of which lie in $OA$, $OB$, $OC$ respectively, and for which
\[
\frac{\sin OX}{\sin OA} = \frac{a}{\theta},\quad \frac{\sin OY}{\sin OB} = \frac{b}{\theta},\quad \frac{\sin OZ}{\sin OC} = \frac{c}{\theta}.
\]
I take the arcs $BC$, $CA$, $AB = \alpha$, $\beta$, $\gamma$ respectively, $\alpha + \beta + \gamma = 2\pi$; and the required arcs $OA$, $OB$, $OC$ are taken to be $\xi$, $\eta$, $\zeta$ respectively; these are connected by the relation
\[
\sin\alpha\cos\xi + \sin\beta\cos\eta + \sin\gamma\cos\zeta = 0,
\]
to obtain which, observe that from the triangles $OAB$, $OAC$, we have
\[
\cos A = \frac{\cos\eta - \cos\xi\cos\gamma}{\sin\xi\sin\gamma} = -\frac{\cos\zeta - \cos\xi\cos\beta}{\sin\xi\sin\beta},
\]
that is,
\[
\sin\beta(\cos\eta - \cos\xi\cos\gamma) + \sin\gamma(\cos\zeta - \cos\xi\cos\beta) = 0,
\]
which, with $\sin\alpha = -\sin(\beta + \gamma)$, gives the required relation. We have
\[
\sin OX = \frac{a}{\theta}\sin\xi,\quad \sin OY = \frac{b}{\theta}\sin\eta,\quad \sin OZ = \frac{c}{\theta}\sin\zeta;
\]
and then from the triangles $OBC$, $OCA$, $OAB$, and the quadrantal triangles $OYZ$, $OZX$, $OXY$, we have
\[
\cos BOC = \frac{\cos\alpha - \cos\eta\cos\zeta}{\sin\eta\sin\zeta} = \frac{-\sqrt{\theta^{2} - b^{2}\sin^{2}\eta}\,\sqrt{\theta^{2} - c^{2}\sin^{2}\zeta}}{bc\sin\eta\sin\zeta\,/\,\theta^{2}},\ \&\text{c.},
\]
that is,
\[
bc(\cos\alpha - \cos\eta\cos\zeta) = -\sqrt{\theta^{2} - b^{2}\sin^{2}\eta}\,\sqrt{\theta^{2} - c^{2}\sin^{2}\zeta},
\]
\[
ca(\cos\beta - \cos\zeta\cos\xi) = -\sqrt{\theta^{2} - c^{2}\sin^{2}\zeta}\,\sqrt{\theta^{2} - a^{2}\sin^{2}\xi},
\]
\[
ab(\cos\gamma - \cos\xi\cos\eta) = -\sqrt{\theta^{2} - a^{2}\sin^{2}\xi}\,\sqrt{\theta^{2} - b^{2}\sin^{2}\eta},
\]
which, when rationalized, are quadric equations in $\cos\xi$, $\cos\eta$, $\cos\zeta$. The first equation, in fact, gives
\[
b^{2}c^{2}(\cos\alpha - \cos\eta\cos\zeta)^{2} = (\theta^{2} - b^{2} + b^{2}\cos^{2}\eta)(\theta^{2} - c^{2} + c^{2}\cos^{2}\zeta),
\]
that is,
\[
(\theta^{2} - b^{2})(\theta^{2} - c^{2}) - b^{2}c^{2}\cos^{2}\alpha + (\theta^{2} - b^{2})c^{2}\cos^{2}\zeta + (\theta^{2} - c^{2})b^{2}\cos^{2}\eta + 2b^{2}c^{2}\cos\alpha\cos\eta\cos\zeta = 0,
\]
or, what is the same thing,
\[
-\left(1 - b^{2} - c^{2} - \frac{b^{2}c^{2}\cos^{2}\alpha}{\theta^{2} - b^{2} \cdot \theta^{2} - c^{2}}\right) + \frac{b^{2}}{\theta^{2} - b^{2}}\cos^{2}\eta + \frac{c^{2}}{\theta^{2} - c^{2}}\cos^{2}\zeta + \frac{2b^{2}c^{2}\cos\alpha}{(\theta^{2} - b^{2})(\theta^{2} - c^{2})}\cos\eta\cos\zeta = 0.
\]
Completing the system, we have
\[
-\left(1 - c^{2} - a^{2} - \frac{c^{2}a^{2}\cos^{2}\beta}{\theta^{2} - c^{2} \cdot \theta^{2} - a^{2}}\right) + \frac{c^{2}}{\theta^{2} - c^{2}}\cos^{2}\zeta + \frac{a^{2}}{\theta^{2} - a^{2}}\cos^{2}\xi + \frac{2c^{2}a^{2}\cos\beta}{(\theta^{2} - c^{2})(\theta^{2} - a^{2})}\cos\zeta\cos\xi = 0,
\]
\[
-\left(1 - a^{2} - b^{2} - \frac{a^{2}b^{2}\cos^{2}\gamma}{\theta^{2} - a^{2} \cdot \theta^{2} - b^{2}}\right) + \frac{a^{2}}{\theta^{2} - a^{2}}\cos^{2}\xi + \frac{b^{2}}{\theta^{2} - b^{2}}\cos^{2}\eta + \frac{2a^{2}b^{2}\cos\gamma}{(\theta^{2} - a^{2})(\theta^{2} - b^{2})}\cos\xi\cos\eta = 0,
\]
\noindent\textbf{[p.~512]}\quad and, as above,
\[
\sin\alpha\cos\xi + \sin\beta\cos\eta + \sin\gamma\cos\zeta = 0.
\]
It seems difficult from these equations to eliminate $\xi$, $\eta$, $\zeta$, so as to obtain an equation in $\theta$; but I employ some geometrical considerations.

Taking $\Pi$ as the pole of the circle $ABC$, and drawing $\Pi X$, $\Pi Y$, $\Pi Z$ to meet the circle in $p$, $q$, $r$ respectively, then, if $\alpha''$, $\beta''$, $\gamma''$ are the cosine-inclinations of $\Pi$ to $X$, $Y$, $Z$ respectively, we have
\[
\sin Xp,\ \sin Yq,\ \sin Zr = \alpha'',\ \beta'',\ \gamma''.
\]
From the right parallel triangles $BYq$ and $CZr$, we have
\[
\sin Yq = \sin BY\sin B,
\]
\[
\sin Zr = \sin CZ\sin C,
\]
and, thence,
\[
\frac{\sin Yq}{\sin Zr} = \frac{\sin BY}{\sin CZ} \cdot \frac{\sin OC}{\sin OB};
\]
or, since
\[
BY = OB - OY,\quad CZ = OC - OZ,
\]
and thence
\[
\sin BY = \frac{1}{\theta}\{\sqrt{\theta^{2} - b^{2}\sin^{2}\eta} - b\cos\eta\},
\]
\[
\sin CZ = \frac{1}{\theta}\{\sqrt{\theta^{2} - c^{2}\sin^{2}\zeta} - c\cos\zeta\},
\]
we obtain
\[
\frac{\beta''}{\gamma''} = \frac{\sqrt{\theta^{2} - b^{2}\sin^{2}\eta} - b\cos\eta}{\sqrt{\theta^{2} - c^{2}\sin^{2}\zeta} - c\cos\zeta}.
\]
We have thence
\[
\beta''\sqrt{\theta^{2} - c^{2}\sin^{2}\zeta} - \gamma''\sqrt{\theta^{2} - b^{2}\sin^{2}\eta} = \beta''c\cos\zeta - \gamma''b\cos\eta,
\]
or, squaring and reducing
\[
\beta''^{2}(\theta^{2} - c^{2}) + \gamma''^{2}(\theta^{2} - b^{2}) + 2\beta''\gamma''\left[-\sqrt{\theta^{2} - c^{2}\sin^{2}\zeta}\sqrt{\theta^{2} - b^{2}\sin^{2}\eta} + bc\cos\eta\cos\zeta\right] = 0;
\]
that is,
\[
\beta''^{2}(\theta^{2} - c^{2}) + \gamma''^{2}(\theta^{2} - b^{2}) + 2\beta''\gamma'' \cdot bc\cos\alpha = 0,
\]
and, similarly,
\[
\gamma''^{2}(\theta^{2} - a^{2}) + \alpha''^{2}(\theta^{2} - c^{2}) + 2\gamma''\alpha'' \cdot ca\cos\beta = 0,
\]
\[
\alpha''^{2}(\theta^{2} - b^{2}) + \beta''^{2}(\theta^{2} - a^{2}) + 2\alpha''\beta'' \cdot ab\cos\gamma = 0,
\]
or, what is the same thing,
\[
\frac{\beta''^{2}}{\theta^{2} - b^{2}} + \frac{\gamma''^{2}}{\theta^{2} - c^{2}} = \frac{-2bc\cos\alpha}{\sqrt{(\theta^{2} - b^{2})(\theta^{2} - c^{2})}} \cdot \frac{\beta''\gamma''}{\sqrt{(\theta^{2} - b^{2})(\theta^{2} - c^{2})}},
\]
\[
\frac{\gamma''^{2}}{\theta^{2} - c^{2}} + \frac{\alpha''^{2}}{\theta^{2} - a^{2}} = \frac{-2ca\cos\beta}{\sqrt{(\theta^{2} - c^{2})(\theta^{2} - a^{2})}} \cdot \frac{\gamma''\alpha''}{\sqrt{(\theta^{2} - c^{2})(\theta^{2} - a^{2})}},
\]
\[
\frac{\alpha''^{2}}{\theta^{2} - a^{2}} + \frac{\beta''^{2}}{\theta^{2} - b^{2}} = \frac{-2ab\cos\gamma}{\sqrt{(\theta^{2} - a^{2})(\theta^{2} - b^{2})}} \cdot \frac{\alpha''\beta''}{\sqrt{(\theta^{2} - a^{2})(\theta^{2} - b^{2})}}.
\]

\noindent\textbf{[p.~513]}\quad writing
\[
\alpha'',\ \beta'',\ \gamma'' = X'\sqrt{\theta^{2} - a^{2}},\ Y'\sqrt{\theta^{2} - b^{2}},\ Z'\sqrt{\theta^{2} - c^{2}},
\]
and
\[
\frac{bc\cos\alpha}{\sqrt{(\theta^{2} - b^{2})(\theta^{2} - c^{2})}} = f,\quad \frac{ca\cos\beta}{\sqrt{(\theta^{2} - c^{2})(\theta^{2} - a^{2})}} = g,\quad \frac{ab\cos\gamma}{\sqrt{(\theta^{2} - a^{2})(\theta^{2} - b^{2})}} = h,
\]
the equations are
\[
Y'^{2} + Z'^{2} - 2fY'Z' = 0,
\]
\[
Z'^{2} + X'^{2} - 2gZ'X' = 0,
\]
\[
X'^{2} + Y'^{2} - 2hX'Y' = 0.
\]
Writing the last two under the form
\[
Z'^{2} - 2gZ'X' + X'^{2} = 0,
\]
\[
X'^{2} - 2hX'Y' + Y'^{2} = 0,
\]
and eliminating $Z'$, we have
\[
-4(1 - g^{2})(1 - h^{2})Y'^{2}X'^{2} + (Y'^{2} + Z'^{2} - 2ghY'Z')^{2} = 0,
\]
which, in virtue of the first equation, is
\[
-4(1 - g^{2})(1 - h^{2})Y'^{2}Z'^{2} + 4(gh - f)^{2}Y'^{2}Z'^{2} = 0,
\]
that is,
\[
(1 - g^{2})(1 - h^{2}) - (gh - f)^{2} = 0;
\]
or, what is the same thing,
\[
1 - f^{2} - g^{2} - h^{2} + 2fgh = 0.
\]
I remark that we may write
\[
gh - f = \sqrt{(1 - g^{2})}\sqrt{(1 - h^{2})},
\]
\[
hf - g = \sqrt{(1 - h^{2})}\sqrt{(1 - f^{2})},
\]
\[
fg - h = \sqrt{(1 - f^{2})}\sqrt{(1 - g^{2})},
\]
the signs on the right-hand side being either all $+$, or else one $+$ and two $-$, so that the product is $+$. In fact, multiplying the assumed equations, we have
\[
f^{2}g^{2}h^{2} - fgh(f^{2} + g^{2} + h^{2}) + f^{2}g^{2} + h^{2}f^{2} + f^{2}g^{2} - fgh = 1 - f^{2} - g^{2} - h^{2} + f^{2}h^{2} + h^{2}f^{2} + f^{2}g^{2} - f^{2}g^{2}h^{2},
\]
that is,
\[
1 - f^{2} - g^{2} - h^{2} + fgh(1 + f^{2} + g^{2} + h^{2}) - 2f^{2}g^{2}h^{2} = 0,
\]
or,
\[
(1 - f^{2} - g^{2} - h^{2} + 2fgh)(1 - fgh) = 0,
\]
which is right; but with a different combination of signs the result would not have been obtained.

Substituting for $f$, $g$, $h$ their values, we have
\[
(a^{2} - \theta^{2})(b^{2} - \theta^{2})(c^{2} - \theta^{2}) - b^{2}c^{2}(a^{2} - \theta^{2})\cos^{2}\alpha - c^{2}a^{2}(b^{2} - \theta^{2})\cos^{2}\beta
\]
\[
- a^{2}b^{2}(c^{2} - \theta^{2})\cos^{2}\gamma + 2a^{2}b^{2}c^{2}\cos\alpha\cos\beta\cos\gamma = 0,
\]
\noindent\textbf{[p.~514]}\quad where the term independent of $\theta$ is
\[
a^{2}b^{2}c^{2}(1 - \cos^{2}\alpha - \cos^{2}\beta - \cos^{2}\gamma + 2\cos\alpha\cos\beta\cos\gamma),
\]
which is $= 0$ in virtue of $\alpha + \beta + \gamma = 2\pi$. We have, therefore, for $\theta^{2}$ the quadric equation
\[
b^{2}c^{2}\sin^{2}\alpha + c^{2}a^{2}\sin^{2}\beta + a^{2}b^{2}\sin^{2}\gamma - (a^{2} + b^{2} + c^{2})\theta^{2} + \theta^{4} = 0,
\]
giving for $\theta^{2}$ the two real positive values
\[
\theta^{2} = \tfrac{1}{2}\{a^{2} + b^{2} + c^{2} + \sqrt{\Omega}\},
\]
where
\[
\Omega = (a^{2} + b^{2} + c^{2})^{2} - 4(b^{2}c^{2}\sin^{2}\alpha + c^{2}a^{2}\sin^{2}\beta + a^{2}b^{2}\sin^{2}\gamma)
\]
\[
= a^{4} + b^{4} + c^{4} + 2b^{2}c^{2}\cos 2\alpha + 2c^{2}a^{2}\cos 2\beta + 2a^{2}b^{2}\cos 2\gamma
\]
\[
= (a^{2} + b^{2}\cos 2\gamma + c^{2}\cos 2\beta)^{2} + (b^{2}\sin 2\gamma - c^{2}\sin 2\beta)^{2}.
\]
I write now
\[
\frac{a\cos\xi}{\sqrt{(a^{2} - \theta^{2})}} = X,\quad \frac{b\cos\eta}{\sqrt{(b^{2} - \theta^{2})}} = Y,\quad \frac{c\cos\zeta}{\sqrt{(c^{2} - \theta^{2})}} = Z,
\]
and also
\[
\frac{\sqrt{a^{2} - \theta^{2}}}{a}\sin\alpha = A,\quad \frac{\sqrt{b^{2} - \theta^{2}}}{b}\sin\beta = B,\quad \frac{\sqrt{c^{2} - \theta^{2}}}{c}\sin\gamma = C.
\]
The equations for $\cos\xi$, $\cos\eta$, $\cos\zeta$ become
\[
Y^{2} + Z^{2} - 2fYZ - (1 - f^{2}) = 0,
\]
\[
Z^{2} + X^{2} - 2gZX - (1 - g^{2}) = 0,
\]
\[
X^{2} + Y^{2} - 2hXY - (1 - h^{2}) = 0,
\]
and
\[
AX + BY + CZ = 0,
\]
in virtue of the relation between $f$, $g$, $h$. The first three equations are satisfied by a two-fold relation between $X$, $Y$, $Z$; viz.\ treating these as coordinates, the equations represent three quadric cylinders having a common conic.

To prove this, I write
\[
1 - f^{2},\ 1 - g^{2},\ 1 - h^{2},\ gh - f,\ hf - g,\ fg - h = \mathfrak{a},\ \mathfrak{b},\ \mathfrak{c},\ \mathfrak{f},\ \mathfrak{g},\ \mathfrak{h}.
\]
We have, as usual,
\[
\mathfrak{b}\mathfrak{c} - \mathfrak{f}^{2},\ \mathfrak{c}\mathfrak{a} - \mathfrak{g}^{2},\ \mathfrak{a}\mathfrak{b} - \mathfrak{h}^{2},\ \mathfrak{g}\mathfrak{h} - \mathfrak{a}\mathfrak{f},\ \mathfrak{h}\mathfrak{f} - \mathfrak{b}\mathfrak{g},\ \mathfrak{f}\mathfrak{g} - \mathfrak{c}\mathfrak{h},\ \text{each} = 0:
\]
the equations
\[
\mathfrak{a}X + \mathfrak{h}Y + \mathfrak{g}Z = 0,\quad \mathfrak{h}X + \mathfrak{b}Y + \mathfrak{f}Z = 0,\quad \mathfrak{g}X + \mathfrak{f}Y + \mathfrak{c}Z = 0,
\]
represent each of them one and the same plane, which I say is that of the conic in question.

\noindent\textbf{[p.~515]}\quad The three given equations are
\[
Y^{2} + Z^{2} - 2fYZ - \mathfrak{a} = 0,
\]
\[
Z^{2} + X^{2} - 2gZX - \mathfrak{b} = 0,
\]
\[
X^{2} + Y^{2} - 2hXY - \mathfrak{c} = 0,
\]
say these are $U = 0$, $V = 0$, $W = 0$; it is to be shown that $cV - bW$, $aW - cU$, $bU - aV$, each contain the linear factor in question. We have
\[
\mathfrak{c}V - \mathfrak{b}W = (\mathfrak{c} - \mathfrak{b})X^{2} - \mathfrak{b}Y^{2} + \mathfrak{c}Z^{2} - 2\mathfrak{c}gZX + 2\mathfrak{b}hXY;
\]
or, what is the same thing,
\[
\mathfrak{a}(\mathfrak{c}V - \mathfrak{b}W) = \mathfrak{a}(\mathfrak{c} - \mathfrak{b})X^{2} - \mathfrak{h}^{2}Y^{2} + \mathfrak{g}^{2}Z^{2} - 2\mathfrak{g}\mathfrak{g}_{*}ZX + 2\mathfrak{h}\mathfrak{h}_{*}XY.
\]
Assuming this
\[
= (\mathfrak{a}X + \mathfrak{h}Y + \mathfrak{g}Z)(\lambda X - \mathfrak{h}Y + \mathfrak{g}Z),
\]
we have
\[
\mathfrak{a}\lambda = \mathfrak{a}(\mathfrak{c} - \mathfrak{b}),
\]
\[
\mathfrak{g}(\mathfrak{a} + \lambda) = -2\mathfrak{g}\mathfrak{g}_{*},
\]
\[
\mathfrak{h}(-\mathfrak{a} + \lambda) = 2\mathfrak{h}\mathfrak{h}_{*},
\]
that is,
\[
\lambda = \mathfrak{c} - \mathfrak{b},\quad \mathfrak{a} + \lambda = -2\mathfrak{g}_{*},\quad -\mathfrak{a} + \lambda = 2\mathfrak{h}_{*};
\]
but $\lambda = \mathfrak{c} - \mathfrak{b}$, $= -\mathfrak{h}^{2} + \mathfrak{g}^{2}$, and the other two equations are $\mathfrak{a} + \mathfrak{c} - \mathfrak{b} + 2\mathfrak{g}_{*} = 0$, $\mathfrak{a} + \mathfrak{b} - \mathfrak{c} + 2\mathfrak{h}_{*} = 0$, which are identically true.

The values of $X$, $Y$, $Z$ are thus determined as the coordinates of the intersection of the conic with the plane $AX + BY + CZ = 0$; or, what is the same thing, of the line
\[
AX + BY + CZ = 0,
\]
\[
\mathfrak{a}X + \mathfrak{h}Y + \mathfrak{g}Z = 0,
\]
with any one of the three cylinders.

We may, however, complete the analytical solution in a different manner as follows. Assuming as above $\sqrt{(\mathfrak{b}\mathfrak{c})} = \mathfrak{f}$, $\sqrt{(\mathfrak{c}\mathfrak{a})} = \mathfrak{g}$, $\sqrt{(\mathfrak{a}\mathfrak{b})} = \mathfrak{h}$, and thence $\mathfrak{h}\sqrt{(\mathfrak{c})} - \mathfrak{g}\sqrt{(\mathfrak{b})} = 0$, we obtain from the second and the third equations
\[
Y = hX + \sqrt{(\mathfrak{c})}\sqrt{(1 - X^{2})},\quad Z = gX - \sqrt{(\mathfrak{b})}\sqrt{(1 - X^{2})},
\]
(the signs are one $+$ the other $-$, in order that this may consist with the equation $\mathfrak{a}X + \mathfrak{h}Y + \mathfrak{g}Z = 0$). Substituting in $AX + BY + CZ = 0$, we have
\[
(A + Bh + Cg)X + \{B\sqrt{(\mathfrak{c})} - C\sqrt{(\mathfrak{b})}\}\sqrt{(1 - X^{2})} = 0,
\]
that is,
\[
(A + Bh + Cg)^{2}X^{2} - (B^{2}\mathfrak{c} + C^{2}\mathfrak{b} - 2BC\mathfrak{f})(1 - X^{2}) = 0,
\]
\noindent\textbf{[p.~516]}\quad or say
\[
(A + Bh + Cg)^{2}X^{2} + \{B^{2}(1 - h^{2}) + C^{2}(1 - g^{2}) - 2BC(gh - f)\}(X^{2} - 1) = 0,
\]
that is,
\[
(A^{2} + B^{2} + C^{2} + 2BCf + 2CAg + 2ABh)X^{2} = [B^{2} + C^{2} + 2BCf - (Bh + Cg)^{2}],
\]
or writing
\[
A^{2} + B^{2} + C^{2} + 2BCf + 2CAg + 2ABh = \Delta,
\]
say we have
\[
\Delta X^{2} = B^{2} + C^{2} + 2BCf - (Bh + Cg)^{2},
\]
\[
\Delta Y^{2} = C^{2} + A^{2} + 2CAg - (Cf + Ah)^{2},
\]
\[
\Delta Z^{2} = A^{2} + B^{2} + 2ABh - (Ag + Bf)^{2}.
\]

Now attending to the values of $A$, $B$, $C$, $f$, $g$, $h$, we have
\[
BCf,\ CAg,\ ABh = \sin\beta\sin\gamma\cos\alpha,\ \sin\gamma\sin\alpha\cos\beta,\ \sin\alpha\sin\beta\cos\gamma,
\]
and thence
\[
\Delta = \sin^{2}\alpha\left(1 - \frac{\theta^{2}}{a^{2}}\right) + \sin^{2}\beta\left(1 - \frac{\theta^{2}}{b^{2}}\right) + \sin^{2}\gamma\left(1 - \frac{\theta^{2}}{c^{2}}\right)
\]
\[
+ 2(\sin\beta\sin\gamma\cos\alpha + \sin\gamma\sin\alpha\cos\beta + \sin\alpha\sin\beta\cos\gamma);
\]
in virtue of $\alpha + \beta + \gamma = 2\pi$, the last term is
\[
= 2(\cos\alpha\cos\beta\cos\gamma - 1),
\]
whence
\[
\Delta = -\theta^{2}\left(\frac{\sin^{2}\alpha}{a^{2}} + \frac{\sin^{2}\beta}{b^{2}} + \frac{\sin^{2}\gamma}{c^{2}}\right),\ \text{say this is}\ = -\theta^{2}\Lambda.
\]
Moreover $Bh + Cg = \dfrac{a\sin\alpha}{\sqrt{(a^{2} - \theta^{2})}}$, whence the value of $\Delta X^{2}$ is
\[
= \sin^{2}\beta\left(1 - \frac{\theta^{2}}{b^{2}}\right) + \sin^{2}\gamma\left(1 - \frac{\theta^{2}}{c^{2}}\right) + 2\sin\beta\sin\gamma\cos\alpha - \left(1 - \frac{\theta^{2}}{a^{2}}\right)^{-1}\sin^{2}\alpha.
\]

Here the constant term is
\[
= \sin^{2}\beta + \sin^{2}\gamma + 2\sin\beta\sin\gamma\cos\alpha,
\]
that is,
\[
= 1 - (1 - \sin^{2}\beta)(1 - \sin^{2}\gamma) + \sin^{2}\beta\sin^{2}\gamma + 2\sin\beta\sin\gamma\cos\alpha
\]
\[
= 1 - \cos^{2}\beta\cos^{2}\gamma - \cos^{2}\alpha + (\cos\alpha + \sin\beta\sin\gamma)^{2}
\]
\[
= 1 - \cos^{2}\alpha, = \sin^{2}\alpha,
\]
or the whole is
\[
\sin^{2}\alpha\left(1 - \frac{a^{2}}{a^{2} - \theta^{2}}\right) - \theta^{2}\left(\frac{\sin^{2}\beta}{b^{2}} + \frac{\sin^{2}\gamma}{c^{2}}\right),
\]
\noindent\textbf{[p.~517]}\quad which is
\[
= -\theta^{2}\left(\frac{\sin^{2}\alpha}{a^{2} - \theta^{2}} + \frac{\sin^{2}\beta}{b^{2}} + \frac{\sin^{2}\gamma}{c^{2}}\right),
\]
so that we have
\[
\Lambda X^{2} = \left(\frac{\sin^{2}\alpha}{a^{2} - \theta^{2}} + \frac{\sin^{2}\beta}{b^{2}} + \frac{\sin^{2}\gamma}{c^{2}}\right).
\]
Similarly,
\[
\Lambda Y^{2} = \left(\frac{\sin^{2}\alpha}{a^{2}} + \frac{\sin^{2}\beta}{b^{2} - \theta^{2}} + \frac{\sin^{2}\gamma}{c^{2}}\right),
\]
\[
\Lambda Z^{2} = \left(\frac{\sin^{2}\alpha}{a^{2}} + \frac{\sin^{2}\beta}{b^{2}} + \frac{\sin^{2}\gamma}{c^{2} - \theta^{2}}\right);
\]
and hence also
\[
\Lambda(1 - X^{2}) = \frac{-\theta^{2}\sin^{2}\alpha}{a^{2}(a^{2} - \theta^{2})},\quad \Lambda(1 - Y^{2}) = \frac{-\theta^{2}\sin^{2}\beta}{b^{2}(b^{2} - \theta^{2})},\quad \Lambda(1 - Z^{2}) = \frac{-\theta^{2}\sin^{2}\gamma}{c^{2}(c^{2} - \theta^{2})},
\]
where
\[
\Lambda = \frac{\sin^{2}\alpha}{a^{2}} + \frac{\sin^{2}\beta}{b^{2}} + \frac{\sin^{2}\gamma}{c^{2}}.
\]

The equation in $X$ is
\[
\Lambda\left(1 - \frac{a^{2}\cos^{2}\xi}{a^{2} - \theta^{2}}\right) = \frac{-\theta^{2}\sin^{2}\alpha}{a^{2}(a^{2} - \theta^{2})},
\]
that is,
\[
\Lambda(a^{2}\sin^{2}\xi - \theta^{2}) = \frac{-\theta^{2}\sin^{2}\alpha}{a^{2}},
\]
or
\[
a^{2}\sin^{2}\xi = \frac{\theta^{2}}{\Lambda}\left(\Lambda - \frac{\sin^{2}\alpha}{a^{2}}\right),
\]
and the like for $\eta$, $\zeta$. Writing for greater convenience $\dfrac{\sin^{2}\alpha}{a^{2}}$, $\dfrac{\sin^{2}\beta}{b^{2}}$, $\dfrac{\sin^{2}\gamma}{c^{2}} = p$, $q$, $r$, then $\Lambda = p + q + r$, and we have
\[
\sin^{2}\xi = \frac{\theta^{2}}{a^{2}} \cdot \frac{q + r}{p + q + r},\quad \sin^{2}\eta = \frac{\theta^{2}}{b^{2}} \cdot \frac{r + p}{p + q + r},\quad \sin^{2}\zeta = \frac{\theta^{2}}{c^{2}} \cdot \frac{p + q}{p + q + r},
\]
(whence also $a^{2}\sin^{2}\xi + b^{2}\sin^{2}\eta + c^{2}\sin^{2}\zeta = 2\theta^{2}$: as a simple verification, observe that, if the projection is rectangular, the axes being all equally inclined to the plane of projection, then $\xi = \eta = \zeta = 90^{\circ}$, $a = b = c = \theta\sin s$, and the equation is $3\sin^{2}s = 2$; $s$, $s$ are here the sides of an isosceles quadrantal triangle, the included angle being $120^{\circ}$, that is, we have $\cos 120^{\circ}$ $(= -\tfrac{1}{2}) = -\cot^{2}s$, that is, $\cot^{2}s = \tfrac{1}{2}$, or $\sin^{2}s = \tfrac{2}{3}$, which is right).

I remark, that a geometrical solution may be obtained upon very different principles. We have on a sphere the trirectangular triangle $XYZ$, which by parallel lines is projected into $ABC$. Every great circle of the sphere is projected into an ellipse having double \noindent\textbf{[p.~518]}\quad contact at the extremities of a diameter with the ellipse which is the apparent contour of the sphere. Moreover, if the arc of great circle $XY$ is a quadrant, then the radius through $X$ and the tangent at $Y$ are parallel to each other, whence, if $\Omega$ be the projection of the centre, and $AB$ the projection of the arc $XY$, then in the projection the line $\Omega A$ and the tangent at $B$ are parallel to each other. It is now easy to derive a construction: with centre $\Omega$, and conjugate semi-axes $(\Omega B, \Omega C)$, $(\Omega C, \Omega A)$, $(\Omega A, \Omega B)$ respectively, describe three ellipses; and find a concentric ellipse having double contact with each of these (there are in fact two such ellipses, one touching the three ellipses internally, and giving an imaginary solution; the other touching them externally, which is the ellipse intended). Drawing then through the ellipse a right cylinder (there are two such cylinders, but only one of them is real), and inscribing in it a sphere, and projecting on to the surface of the sphere by lines parallel to the axis of the cylinder, the three ellipses are projected into three great circles cutting at right angles, or, say, the elliptic arcs $BC$, $CA$, $AB$ are projected into the trirectangular triangle $XYZ$.

\newpage

%% ========================================================================
%% PAPER [615] -- On the Conic Torus
%%                Printed pages 519--521 (PDF pp. 541--543)
%% ========================================================================
\section*{\textbf{615.} On the Conic Torus}
\addcontentsline{toc}{section}{\textbf{615.} On the Conic Torus (pp.~519--521)}

\textit{[From the Quarterly Journal of Pure and Applied Mathematics, vol.~xiii. (1875), pp.~127--129.]}

\bigskip

\noindent\textbf{[p.~519]}\quad The equation $p + \sqrt{(qr)} + \sqrt{(st)} = 0$, where $p$, $q$, $r$, $s$, $t$ are linear functions of the coordinates $(x, y, z, w)$, and as such are connected by a linear relation, belongs to a quartic surface having a nodal conic $(p = 0, qr - st = 0)$; and four nodes (conical points), viz.\ these are the intersections of the line $q = 0$, $r = 0$ with the quadric surface $p^{2} - qr - st = 0$, and of the line $r = 0$, $s = 0$ with the same surface. The quartic surface has also four tropes (planes which touch the surface along a conic); viz.\ these are the planes $q = 0$, $r = 0$, $s = 0$, $t = 0$, the conic of contact or tropal conic in each plane being the intersection of the plane with the before-mentioned quadric surface $p^{2} - qr - st = 0$. The planes $q = 0$, $r = 0$, and also the conies in these planes pass through two of the nodes, say $A$, $C$; and the planes $s = 0$, $t = 0$, and also the conies in these planes pass through the remaining two nodes, say $B$, $D$; so that the relations of the surface are as is shown in fig.~1. It is to be added that $AB$, $BC$, $CD$, $DA$ (but not $AC$ or $BD$) are lines on the surface.

The planes $q = 0$, $r = 0$, which contain the tropal conies through $A$, $C$, are in general distinct from the planes $ABC$, $ADC$ which contain the line-pairs $BA$, $BC$ and $DA$, $DC$ respectively: and so also the planes $s = 0$, $t = 0$, which contain the tropal conies through $B$, $D$, are in general distinct from the planes $ABD$, $CBD$ which contain the line-pairs $AB$, $AD$ and $CB$, $CD$ respectively.

If, however, the identical linear relation contain only $p$, $s$, $t$, then the planes $q = 0$, $r = 0$ will be the planes $ABC$, $ADC$ respectively: and the tropal conies in these planes will consequently be the line-pairs $BA$, $BC$, and $DA$, $DC$ respectively. But the planes $s = 0$, $t = 0$ will continue to be distinct from the planes $ABD$, $CBD$: and the tropal conies in the planes $s = 0$, $t = 0$ will remain proper conies.

\noindent\textbf{[p.~520]}\quad A surface of the last-mentioned form is
\[
mz + \sqrt{(xy)} + \sqrt{(w^{2} - z^{2})} = 0,
\]
viz.\ this has the nodal conic $z = 0$, $xy - w^{2} = 0$, the nodes $\{x = 0, y = 0, (m^{2} + 1)z^{2} - w^{2} = 0\}$, and $(z = 0, w = 0, x = 0)$, $(z = 0, w = 0, y = 0)$, and the tropes $x = 0$, $y = 0$, $z + w = 0$, $z - w = 0$; but the planes $z + w = 0$ and $z - w = 0$ are ordinary tropal planes each touching the surface in a proper conic; the planes $x = 0$, $y = 0$ are special planes each touching along a line-pair.

\begin{center}
\textit{[Fig.~1: Tetrahedron $ABCD$ with the marked planes and edges as described.]}
\end{center}

The equation in question, writing therein $w = 1$ and $x + iy$, $x - iy$ in place of $(x, y)$ respectively, is
\[
\sqrt{(x^{2} + y^{2})} + (mz)^{2} = 1 - z^{2},
\]
which is derived from
\[
(x + mz)^{2} = 1 - z^{2},
\]
by the change of $x$ into $\sqrt{(x^{2} + y^{2})}$; and the surface is consequently the torus generated by the rotation of the conic $(x + mz)^{2} = 1 - z^{2}$ about its diameter. Or, what is the same thing, the surface
\[
mz + \sqrt{(xy)} + \sqrt{(w^{2} - z^{2})} = 0,
\]
regarding therein $(x, y)$ as circular coordinates and $w$ as being $= 1$, is a torus. The rational equation is $U = 0$, where we have
\[
U = \{(m^{2} + 1)z^{2} - w^{2} + xy\}^{2} - 4m^{2}z^{2}xy
\]
\[
= \{xy + (1 - m^{2})z^{2} - w^{2}\}^{2} + 4m^{2}z^{2}(z^{2} - w^{2})
\]
\[
= x^{2}y^{2} + (m^{2} + 1)^{2}z^{4} + w^{4} + (2 - 2m^{2})z^{2}xy - (2 + 2m^{2})z^{2}w^{2} - 2xyw^{2}.
\]
I find that the Hessian $H$ of this function $U$ contains the factor $xy + (1 - m^{2})z^{2} - w^{2}$, viz.\ that we have
\[
H = \{xy + (1 - m^{2})z^{2} - w^{2}\}H',
\]
\noindent\textbf{[p.~521]}\quad where
\[
H' = x^{3}y^{3}(1 - m^{2})
\]
\[
+ x^{2}y^{2}\left[(3 + 8m^{2} + m^{4})z^{2} + (-3 + m^{2})w^{2}\right]
\]
\[
+ xy\left\{(3 + 11m^{2} + 9m^{4} + m^{6})z^{4} + (-6 - 12m^{2} + 6m^{4})z^{2}w^{2} + (3 + m^{2})w^{4}\right\}
\]
\[
+ (1 + m^{2})\{(1 - m^{2})z^{2} - w^{2}\}\{(1 + m^{2})z^{2} - w^{2}\}^{2},
\]
giving without much difficulty
\[
H' = z^{6}(1 + m^{2})^{3}(1 - m^{2})
\]
\[
+ 2z^{4}\left[(1 + 4m^{2} + m^{4})xy - (1 - m^{4})w^{2}\right](1 + m^{2})
\]
\[
+ z^{2}(xy - w^{2})\left[(1 + 12m^{2} - m^{4})xy - (1 - m^{4})w^{2}\right]
\]
\[
+ \left[(1 - m^{2})xy - (1 + m^{2})w^{2}\right]U;
\]
say this is
\[
= z^{2}H'' + [(1 - m^{2})xy - (1 + m^{2})w^{2}]U,
\]
where
\[
H'' = z^{4}(1 + m^{2})^{2}(1 - m^{2})
\]
\[
+ 2z^{2}(1 + m^{2})\left[(1 + 4m^{2} + m^{4})xy - (1 - m^{4})w^{2}\right]
\]
\[
+ (xy - w^{2})\left[(1 + 12m^{2} - m^{4})xy - (1 - m^{4})w^{2}\right],
\]
or, what is the same thing,
\[
H'' = x^{2}y^{2}(1 + 12m^{2} - m^{4})
\]
\[
+ 2xy\left[(1 + 4m^{2} + m^{4})(1 + m^{2})z^{2} + (-1 + 6m^{2})w^{2}\right]
\]
\[
+ (1 - m^{4})\{(1 + m^{2})z^{2} - w^{2}\}^{2}.
\]

It consequently appears that the complete spinode curve or intersection of the quartic surface and its Hessian, being of order $4 \times 8$, $= 32$, breaks up into
\[
U = 0,\quad xy + (1 - m^{2})z^{2} - w^{2} = 0,
\]
that is,
\[
\begin{array}{ll}
\text{conic}\ z = 0,\ xy - w^{2} = 0\ \text{twice,} & \text{order}\ 4 \\
\text{conic}\ z + w = 0,\ xy - m^{2}w^{2} = 0, & \quad\quad\ \ 2 \\
\text{conic}\ z - w = 0,\ xy - m^{2}w^{2} = 0, & \quad\quad\ \ 2
\end{array}
\]
and
\[
U = 0,\quad z^{2}H'' = 0,
\]
that is,
\[
\begin{array}{ll}
U = 0,\ z^{2} = 0,\ \text{conic}\ z = 0,\ xy - w^{2} = 0\ \text{four times,} & \text{order}\ 8 \\
\text{proper spinode curve}\ U = 0,\ H'' = 0, & \quad\quad\ \ 16 \\
\end{array}
\]
\[
\overline{\phantom{xxxxxxxxxxxxxxxxxxxxxxxxxxxxxxxxxxxxxxxxxxxxxxxxxxx}}\ 32;
\]
viz.\ the intersection is made up of the conic $z = 0$, $xy - w^{2} = 0$ six times, the conies $z \pm w = 0$, $xy - m^{2}w^{2} = 0$ each twice, and the proper spinode curve of the order 16.

\newpage

%% ========================================================================
%% PAPER [616] -- A Geometrical Illustration of the Cubic Transformation in Elliptic Functions
%%                Printed pages 522--526 (PDF pp. 544--548)
%% ========================================================================
\section*{\textbf{616.} A Geometrical Illustration of the Cubic Transformation in Elliptic Functions}
\addcontentsline{toc}{section}{\textbf{616.} A Geometrical Illustration of the Cubic Transformation in Elliptic Functions (pp.~522--526)}

\textit{[From the Quarterly Journal of Pure and Applied Mathematics, vol.~xiii. (1875), pp.~211--216.]}

\bigskip

\noindent\textbf{[p.~522]}\quad Consider the cubic curve
\[
x^{3} + y^{3} + z^{3} + 6lxyz = 0.
\]

If through one of the inflexions $z = 0$, $x + y = 0$, we draw an arbitrary line $z = u(x + y)$, we have at the other intersections of this line with the curve
\[
u\{u^{2}(x + y)^{2} + 6lxy\} + x^{2} - xy + y^{2} = 0;
\]
that is,
\[
(u^{3} + 1)(x^{2} + y^{2}) + 2xy(u^{3} + 3lu - \tfrac{1}{2}) = 0;
\]
and from this equation it appears that the ratio $x : y$ is given as a function involving the square root of
\[
(u^{3} + 3lu - \tfrac{1}{2})^{2} - (u^{3} + 1)^{2},
\]
which, rejecting a factor 3, is
\[
= (2u^{3} + 3lu + \tfrac{1}{2})(lu - \tfrac{1}{2}).
\]
It may be noticed that $lu - \tfrac{1}{2} = 0$ gives the value of $u$, which in the equation $z = u(x + y)$ belongs to the tangent at the inflexion; and $2u^{3} + 3lu + \tfrac{1}{2} = 0$ gives the values which belong to the three tangents from the inflexion.

It thus appears that the coordinates $x$, $y$, $z$ of any point of the curve can be expressed as proportional to functions of $u$ involving the radical
\[
\sqrt{\{(lu - \tfrac{1}{2})(2u^{3} + 3lu + \tfrac{1}{2})\}},
\]
and the theory of the curve is connected with that of a quasi-elliptic integral depending on this radical.

\noindent\textbf{[p.~523]}\quad Taking $\omega$ an imaginary cube root of unity, write
\[
\omega x + \omega^{2}y - 2lz = x',
\]
\[
\omega^{2}x + \omega y - 2lz = y',
\]
\[
x + y - 2lz = z';
\]
then we have
\[
x'y'z' = x^{3} + y^{3} - 8l^{3}z^{3} + 6lx^{2}yz = x^{3} + y^{3} + z^{3} + 6lxyz - (1 + 8l^{3})z^{3}.
\]
Also
\[
-6lz = x' + y' + z',\quad z^{3} = \frac{-1}{216l^{3}}(x' + y' + z')^{3},
\]
whence
\[
x'^{3} + y'^{3} + z'^{3} - \frac{216l^{3}}{1 + 8l^{3}}x'y'z' = \frac{216l^{3}}{1 + 8l^{3}}(x'^{3} + y'^{3} + z'^{3} + 6lx'y'z'),
\]
so that, putting
\[
6m = \frac{-216l^{3}}{1 + 8l^{3}},
\]
or, what is the same thing,
\[
8l^{3}m^{3} + l^{3} + m^{3} = 0,
\]
the equation of the curve is
\[
(x' + y' + z')^{3} + 216m^{3}x'y'z' = 0;
\]
and if we write
\[
x' : y' : z' = X^{3} : Y^{3} : Z^{3},
\]
then the original curve is transformed into
\[
(X^{3} + Y^{3} + Z^{3})^{3} + 216m^{3}X^{3}Y^{3}Z^{3} = 0,
\]
a curve of the ninth order breaking up into three cubic curves, one of which is
\[
X^{3} + Y^{3} + Z^{3} + 6mXYZ = 0,
\]
and for the other two we write herein $m\omega$ and $m\omega^{2}$ respectively in place of $m$. Attending only to the first curve, we have
\[
x^{3} + y^{3} + z^{3} + 6lxyz = 0,
\]
\[
X^{3} + Y^{3} + Z^{3} + 6mXYZ = 0,
\]
as corresponding curves, the corresponding points being connected by the relation
\[
\omega x + \omega^{2}y - 2lz : \omega^{2}x + \omega y - 2lz : x + y - 2lz = X^{3} : Y^{3} : Z^{3},
\]
or, for convenience, we may write
\[
\omega x + \omega^{2}y - 2lz = X^{3},\quad \text{giving}\quad 3x = \omega^{2}X^{3} + \omega Y^{3} + Z^{3},
\]
\[
\omega^{2}x + \omega y - 2lz = Y^{3},\quad \phantom{\text{giving}}\quad 3y = \omega X^{3} + \omega^{2}Y^{3} + Z^{3},
\]
\[
x + y - 2lz = Z^{3},\quad \phantom{\text{giving}}\quad -6lz = X^{3} + Y^{3} + Z^{3}.
\]

\noindent\textbf{[p.~524]}\quad This is a $(1, 3)$ correspondence; viz.\ to a given point on the curve $(m)$, there corresponds one point on $(l)$; but to a given point on $(l)$, three points on $(m)$. As to the first case, this is obvious. As to the second case, if the point $(x, y, z)$ is given, then the corresponding point $(X, Y, Z)$ on the other curve will lie on one of the three lines
\[
Y^{3}(\omega x + \omega^{2}y - 2lz) - X^{3}(\omega^{2}x + \omega y - 2lz) = 0;
\]
each of these intersects the curve $(m)$ in three points: but of the points in the same line it is only one which is a corresponding point of $(x, y, z)$, and the number of the corresponding points is consequently the same as the number of lines, viz.\ it is $= 3$.

We infer that the above equations lead to a cubic transformation of the quasi-elliptic integral
\[
\int du \div \sqrt{\{(lu - \tfrac{1}{2})(2u^{3} + 3lu + \tfrac{1}{2})\}},
\]
into one of the like form
\[
\int dv \div \sqrt{\{(mv - \tfrac{1}{2})(2v^{3} + 3mv + \tfrac{1}{2})\}};
\]
and this is now to be verified.

We have, as before, the line $z = u(x + y)$ meeting the curve $(l)$ in the points
\[
(u^{3} + 1)(x^{2} + y^{2}) + 2xy(u^{3} + 3lu - \tfrac{1}{2}) = 0;
\]
and if similarly through an inflexion of the curve $(m)$ we take the line $Z = v(X + Y)$, this meets the curve in the points
\[
(v^{3} + 1)(X^{2} + Y^{2}) + 2XY(v^{3} + 3mv - \tfrac{1}{2}) = 0.
\]
Then if $(x, y, z)$, $(X, Y, Z)$ are taken to be the corresponding points as above, we can obtain $v$ as a function of $u$. We, in fact, have
\[
-2lu = \frac{-2lz}{x + y} = \frac{X^{3} + Y^{3} + Z^{3}}{-X^{3} - Y^{3} + 2Z^{3}} = \frac{X^{3} + Y^{3} + v^{3}(X + Y)^{3}}{-(X^{3} + Y^{3}) + 2v^{3}(X + Y)^{3}}
\]
\[
= \frac{X^{2} - XY + Y^{2} + v^{3}(X + Y)^{2}}{-X^{2} + XY - Y^{2} + 2v^{3}(X + Y)^{2}}
\]
\[
= \frac{(v^{3} + 1)(X^{2} + Y^{2}) + (2v^{3} - 1)XY}{(2v^{3} - 1)(X^{2} + Y^{2}) + (4v^{3} + 1)XY},
\]
or, since we have
\[
(v^{3} + 1)(X^{2} + Y^{2}) + 2XY(v^{3} + 3mv - \tfrac{1}{2}) = 0,
\]
that is,
\[
X^{2} + Y^{2} : XY = -2v^{3} - 6mv + 1 : v^{3} + 1,
\]
the equation becomes
\[
-2lu = \frac{-6mv(v^{3} + 1)}{(2v^{3} - 1)(-2v^{3} - 6mv + 1) + (4v^{3} + 1)(v^{3} + 1)}
\]
\[
= \frac{-6mv(v^{3} + 1)}{-3v(4mv^{3} - 3v^{2} - 2m)},
\]
\noindent\textbf{[p.~525]}\quad or say,
\[
-lu = m(v^{3} + 1) \div (\square),\quad \text{where the denominator}\ = 4mv^{3} - 3v^{2} - 2m.
\]
This may also be written
\[
-(lu - \tfrac{1}{2}) = 3v^{2}(mv - \tfrac{1}{2}) \div \square.
\]
Proceeding to calculate $2u^{3} + 3lu + \tfrac{1}{2}$, omitting the denominator $(4mv^{3} - 3v^{2} - 2m)^{3}$, this is
\[
\frac{-2}{l^{3}}(v^{3} + 1)^{3} - 3m(v^{3} + 1)(4mv^{3} - 3v^{2} - 2m)^{2} + \tfrac{1}{2}(4mv^{3} - 3v^{2} - 2m)^{3};
\]
or, observing that
\[
m^{3} = \frac{-l^{3}}{1 + 8l^{3}},
\]
that is,
\[
\frac{1}{l^{3}} = \frac{-1 - 8m^{3}}{m^{3}}\quad \text{or}\quad \frac{m^{3}}{l^{3}} = -1 - 8m^{3},
\]
the numerator is
\[
= 2(1 + 8m^{3})(v^{3} + 1)^{3} - 3m(v^{3} + 1)(4mv^{3} - 3v^{2} - 2m)^{2} + \tfrac{1}{2}(4mv^{3} - 3v^{2} - 2m)^{3},
\]
which is found to be identically
\[
\tfrac{1}{2}(2v^{3} + 3mv + \tfrac{1}{2})(v^{3} + 6mv - 2)^{2},
\]
viz.\ we have
\[
2u^{3} + 3lu + \tfrac{1}{2} = (2v^{3} + 3mv + \tfrac{1}{2})(v^{3} + 6mv - 2)^{2} \div (4mv^{3} - 3v^{2} - 2m)^{3}.
\]
and hence
\[
(lu - \tfrac{1}{2})(2u^{3} + 3lu + \tfrac{1}{2}) = -3(mv - \tfrac{1}{2})(2v^{3} + 3mv + \tfrac{1}{2})(v^{3} + 6mv - 2)^{2}v^{2} \div (4mv^{3} - 3v^{2} - 2m)^{4}.
\]
Moreover, we find
\[
ldu = 3mdv \cdot v(v^{3} + 6mv - 2) \div (4mv^{3} - 3v^{2} - 2m)^{2},
\]
and we thence have
\[
\frac{ldu}{\sqrt{\{(lu - \tfrac{1}{2})(2u^{3} + 3lu + \tfrac{1}{2})\}}} = \sqrt{(-3)} \cdot \frac{mdv}{\sqrt{\{(mv - \tfrac{1}{2})(2v^{3} + 3mv + \tfrac{1}{2})\}}},
\]
viz.\ this differential equation corresponds to the integral equation
\[
-lu = m(v^{3} + 1) \div (4mv^{3} - 3v^{2} - 2m),
\]
where $8l^{3}m^{3} + l^{3} + m^{3} = 0$, which corresponds to the modular equation.

It may be remarked that, if $v$ is the same function of $u'$, $l$, in that $u$ is of $v$, $m$, $l$; viz.\ if
\[
-mv = l(u'^{3} + 1) \div (4lu'^{3} - 3u'^{2} - 2m'),
\]
then
\[
\frac{mdv}{\sqrt{\{(mv - \tfrac{1}{2})(2v^{3} + 3mv + \tfrac{1}{2})\}}} = \sqrt{(-3)} \cdot \frac{-ldu'}{\sqrt{\{(lu' - \tfrac{1}{2})(2u'^{3} + 3lu' + \tfrac{1}{2})\}}},
\]
and consequently
\[
\frac{du}{\sqrt{\{(lu - \tfrac{1}{2})(2u^{3} + 3lu + \tfrac{1}{2})\}}} = \frac{-3du'}{\sqrt{\{(lu' - \tfrac{1}{2})(2u'^{3} + 3lu' + \tfrac{1}{2})\}}},
\]
which accords with the general theory of the cubic transformation.

\noindent\textbf{[p.~526]}\quad We may inquire into the relation between the absolute invariants of the two curves. Taking the absolute invariant to be
\[
\Omega = \frac{64S^{3} - T^{2}}{64S^{3}},
\]
where $S$ and $T$ bear the usual significations, we have for the one curve
\[
\Omega = \frac{(1 + 8l^{3})^{3}}{64l^{3}(1 - l^{3})^{3}},
\]
and for the other curve
\[
\Omega' = \frac{(1 + 8m^{3})^{3}}{64m^{3}(1 - m^{3})^{3}},
\]
and, as above, $8l^{3}m^{3} + l^{3} + m^{3} = 0$: writing herein
\[
l^{3} = \frac{-\alpha'}{8\alpha'},\quad m^{3} = \frac{-\beta'}{8\beta'},
\]
the relation between $\alpha'$, $\beta'$ is simply $\alpha' + \beta' = 1$; and the values of $\Omega$, $\Omega'$ are found to be
\[
\Omega = \frac{64\alpha'(1 - \alpha')^{3}}{(1 + 8\alpha')^{3}},\quad \Omega' = \frac{64\beta'(1 - \beta')^{3}}{(1 + 8\beta')^{3}},
\]
viz.\ the required relation is given by the elimination of $\alpha'$, $\beta'$ from these three equations. Or, what is the same thing, writing $\alpha' = \tfrac{1}{2} + \theta$, and therefore $\beta' = \tfrac{1}{2} - \theta$, we have
\[
(5 + 8\theta)^{3}\Omega = 4(1 + 2\theta)(1 - 2\theta)^{3},
\]
\[
(5 - 8\theta)^{3}\Omega' = 4(1 + 2\theta)^{3}(1 - 2\theta),
\]
and the elimination of $\theta$ from these equations gives the required relation between $\Omega$ and $\Omega'$.

It of course follows that, if we have a cubic transformation
\[
\frac{dx}{\sqrt{\{(a, b, c, d, e \mid x, 1)^{4}\}}} = \frac{Cdx}{\sqrt{\{(a', b', c', d', e' \mid x, 1)^{4}\}}},
\]
then the absolute invariants $\Omega$, $\Omega'$ of the two quartic functions are connected by the above relation. I have obtained this result, by reducing the radicals to the standard forms
\[
\sqrt{(1 - x^{2} \cdot 1 - k^{2}x^{2})},\ \sqrt{(1 - x^{2} \cdot 1 - \lambda^{2}x^{2})},
\]
from the known modular equation as represented by the equations
\[
k^{2} = \frac{\alpha^{3}(2 + \alpha)}{1 + 2\alpha},\quad \lambda^{2} = \frac{\alpha(2 + \alpha)^{3}}{(1 + 2\alpha)^{3}};
\]
viz.\ the values of the absolute invariants
\[
\left(= 1 - \frac{27}{2J^{2}}\,,\ 1 - \frac{27J'^{2}}{2J'^{2}}\right)
\]
are
\[
\frac{(k^{4} + 14k^{2} + 1)^{3}}{}\,,\quad \frac{(\lambda^{4} + 14\lambda^{2} + 1)^{3}}{}\,,
\]
but the method of effecting this is by no means obvious.

\newpage

%% ========================================================================
%% PAPER [617] -- On the Scalene Transformation of a Plane Curve
%%                Printed pages 527--528+ (PDF pp. 549--550, continues beyond)
%% ========================================================================
\section*{\textbf{617.} On the Scalene Transformation of a Plane Curve}
\addcontentsline{toc}{section}{\textbf{617.} On the Scalene Transformation of a Plane Curve (pp.~527--528, continues)}

\textit{[From the Quarterly Journal of Pure and Applied Mathematics, vol.~xiii. (1875), pp.~321--328.]}

\bigskip

\noindent\textbf{[p.~527]}\quad The transformation by reciprocal radius vectors can be effected mechanically by Sylvester's Peaucellier-cell. But, employing a more general cell (considered incidentally by him) which may be called the scalene-cell, we have the scalene transformation in question\footnote{The transformation itself, and doubtless many of the results obtained by means of it, are familiar to Prof.~Sylvester; and I abandon all claim to priority.}; viz.\ if, in two curves, $r$, $r'$ are radius vectors belonging to the same angle (or say opposite angles) $\theta$, then the relation between $r$, $r'$ is
\[
rr'(r + r') + (m^{2} - l^{2})r' + (m^{2} - n^{2})r = 0;
\]
or, as this may also be written,
\[
r^{2} + \left(r' - \frac{l^{2} - m^{2}}{r'}\right)r + m^{2} - n^{2} = 0.
\]
The transformation is, it will be seen, an interesting one for its own sake, independently of the remarkably simple mechanical construction, viz.\ the scalene cell is simply a system of 3 pairs of equal rods $PA$, $QA$; $PB$, $QB$; $PC$, $QC$ (fig.~1, p.~528), jointed together at and capable of rotating about the points $P$, $Q$, $A$, $B$, $C$; the three lengths $PA$, $PB$, $PC$ (say these are $= l$, $m$, $n$) are all of them unequal: in the case of any two of them equal, we have Peaucellier or isosceles cell. The effect of the arrangement is that the points $A$, $B$, $C$ are retained in a right line, the distances $BA$, $= r'$, and $BC$, $= r$, being connected by the above-mentioned equation; so that taking $B$ as a fixed point, if the point $A$ describe any given curve, the point $C$ will describe the corresponding or transformed curve.

In the case where the given curve is a right line or a circle, we may through $B$ draw at right angles to the curve the axis $x'Bx$: viz.\ in the case of the circle, \noindent\textbf{[p.~528]}\quad the axis $x'Bx$ passes through its centre; and we measure the angle $\theta$ from this line, viz.\ we write $\angle xBC = \angle x'BA = \theta$.

\begin{center}
\textit{[Fig.~1: The scalene cell with hinge points $P$, $Q$, $A$, $B$, $C$ and rods of lengths $l$, $m$, $n$.]}
\end{center}

Suppose, first, that the locus of $A$ is a right line, or a circle passing through $B$. Its equation is $r' = \dfrac{c}{\cos\theta}$ or $= c\cos\theta$; and we accordingly have for the transformed curve
\[
r^{2} + \left(\frac{c}{\cos\theta} - \frac{l^{2} - m^{2}}{c\cos\theta\,/\,1}\right)r + m^{2} - n^{2} = 0,
\]
or else
\[
r^{2} + \left(c\cos\theta - \frac{l^{2} - m^{2}}{c\cos\theta}\right)r + m^{2} - n^{2} = 0:
\]
viz.\ multiplying in each case by $r\cos\theta$, and then writing $r\cos\theta = x$, $r^{2} = x^{2} + y^{2}$, the equations become
\[
x(x^{2} + y^{2}) + c(x^{2} + y^{2}) - \frac{l^{2} - m^{2}}{c}(x^{2} + y^{2}) + (m^{2} - n^{2})x = 0,
\]
and
\[
x(x^{2} + y^{2}) + cx^{2} - \frac{l^{2} - m^{2}}{c}(x^{2} + y^{2}) + (m^{2} - n^{2})x = 0;
\]
viz.\ in each case the curve is a circular cubic passing through the origin $B$ and having an asymptote parallel to the axis of $y$. The curve is nodal, if $m = n$, viz.\ in this case the origin is a node: or if $c = \sqrt{(l^{2} - n^{2})} + \sqrt{(m^{2} - n^{2})}$.

Suppose next that the locus of $A$ is a circle, centre at a distance $= \gamma$ along $Bx'$ and radius $= h$: we have
\[
r'^{2} - 2\gamma r'\cos\theta + \gamma^{2} - h^{2} = 0,
\]
viz.\ if
\[
\gamma^{2} - h^{2} = -(l^{2} - m^{2}),
\]
or, what is the same thing,
\[
h^{2} + m^{2} = \gamma^{2} + l^{2},
\]
then we have
\[
r' - \frac{l^{2} - m^{2}}{r'} = 2\gamma\cos\theta,
\]
\textit{[Continued in next instalment, pp.~528--534.]}

\end{document}
